\chapter{Problem Context}
\label{chap:problem-context}

\section{Description of the Selected Vietnamese Traffic Scenario}

\fillin{Describe the specific traffic scenario your group chose, for example finding a route between districts in Ho Chi Minh City during rush hour, or planning a delivery route across several wards in Hanoi.}

The scenario used in this project is set in \fillin{city name}, and it covers the area of \fillin{describe the area, for example a set of districts, a university campus, or an industrial zone}. The locations involved and the reason they were chosen are explained in more detail in the dataset section of this report.

\section{Explanation of the Real World Problem Being Addressed}

Vietnamese cities, particularly the large ones such as Ho Chi Minh City and Hanoi, face a combination of high motorbike density, narrow streets, and irregular traffic flow. A driver who only looks at the shortest distance between two points on a map often ends up on a route that is slower in practice because of congestion, one way streets, or poor road quality. The real world problem this project addresses is:

\begin{itemize}[itemsep=2pt]
  \item How to find a route between two locations that is efficient not only in distance but also in travel time.
  \item How to take congestion level and road type into account when more than one path exists between the same two points.
  \item How to plan a trip that must pass through several locations in an order that reduces total time or distance, similar to a delivery run or an errand route.
\end{itemize}

\fillin{Add any details specific to your own scenario, for example a delivery company, a ride hailing service, or a personal commute problem that motivated the choice of this scenario.}

\section{Why Route Optimization Is Useful in This Context}

Route optimization is useful here for several concrete reasons.

\begin{enumerate}[itemsep=2pt]
  \item \textbf{Time saving.} Even a small reduction in travel time per trip becomes significant when multiplied across many trips per day, which matters for delivery services and ride hailing drivers.
  \item \textbf{Fuel and cost saving.} A shorter or less congested route reduces fuel consumption and vehicle wear.
  \item \textbf{Congestion awareness.} Traffic in Vietnamese cities changes quickly during the day. A system that models congestion as part of the cost of a route can recommend a path that is realistic for the current conditions instead of a path that is only correct on paper.
  \item \textbf{Scalability to multiple stops.} Many real tasks are not a single trip from A to B but a sequence of stops, so a method that also optimizes the order of visits adds real value beyond a simple point to point route.
\end{enumerate}

This project applies classic search algorithms studied in the Introduction to Artificial Intelligence course to this practical setting, treating the road network as a graph and the route finding task as a search problem, which is described in detail in the next chapter.
